\section{Introduction}
% the definition of OSS and OSS license
Open source software (OSS) significantly facilitates software development. Developers do not need to reinvent the wheel but only focus on their unique workflows and features. Despite the benefits, misuse of OSS might also induce security issues and legal risks~\cite{licensing2004software,infringement-case,chinacase,infringement-case2,liu2022demystifying}. Many issues arise from OSS licenses, which regulate the conditions under which OSS can be \textit{legally} reused, distributed, and modified. Generally, an OSS license plays the role as a contract signed between the software owner and the user who reuses the code or binary. Once a user of OSS does not conform to the conditions dictated by OSS licenses, license violation occurs and legal risks are induced such as copyright infringement~\cite{infringement-case3}. 

OSS licenses, which are legally binding, protect both the rights of software owners and the freedom of OSS.
Nevertheless, a common issue arises when incorporating third-party OSS accompanied with licenses, i.e., \textit{license incompatibility}, which occurs when multiple OSS licenses exist in the same project and there are conflicts between the conditions they state~\cite{gangadharan2012managing,jss`17-Kapitsaki-SPDX}. For example, ``CANNOT sublicense'' declared by the license of a third-party component but ``CAN sublicense'' declared by the license of the whole project {are incompatible with each other}. In this case, conforming to the project license cannot ensure that the requirements declared by the license of the third-party component are satisfied simultaneously. License incompatibility can also occur among the licenses of multiple third-party components, 
%For example, 
{e.g.,} ``MUST contact author'' from the license of one component but ``CANNOT contact author'' from the license of another component. In this case, at least one license is being violated due to the license incompatibility issue.

License incompatibility represents a serious threat to all stakeholders from a legal perspective. An empirical study on 1,846 GitHub projects shows that 72.91\% of the investigated projects suffered from license incompatibility~\cite{xu2021lidetector}. Despite being problematic, fixing license incompatibility issues requires substantial efforts to either replace the reused code or migrate licenses. Worse still, developers struggle when multiple licenses are involved in the same project~\cite{almeida2017software}. To address this issue, Kapitsaki et al.~\cite{jss`17-Kapitsaki-SPDX,kapitsaki2019modeling} proposed a solution based on the Software Package Data Exchange (SPDX) specification~\cite{spdx}. They designed a directed graph with the compatibility relationships between licenses, based on which
%Based on the graph,
they proposed SPDX Violation Tools (SPDX-VT) to detect license incompatibility and recommend licenses to handle incompatibility issues. Besides, Liu et al.~\cite{liu2019predicting} also proposed Automatic License Prediction (ALP) to prevent license incompatibility issues induced by software changes. 

Despite the progress, there are mainly three problems that limit the application of previous studies. \textbf{First}, previous studies only cover a small predefined set of licenses (e.g., 20 licenses supported by SPDX-VT~\cite{kapitsaki2019modeling} and 25 licenses supported by ALP~\cite{liu2019predicting}). However, there are 489 licenses and versions in SPDX~\cite{spdx}, not to mention potentially a considerable %amount
{number} of license exceptions and custom licenses. \textbf{Second}, previous works only considered the names and versions of some well-known licenses, lacking the capability to comprehend license texts and understand the inherent reasons of license incompatibility, which limits their abilities to resolve incompatibility issues in a flexible and fine-grained way. % when licenses are changed, updated, or customized by developers. 
\textbf{Third}, the solutions to license incompatibility have been mainly focused on recommending an official license as the new license for the whole project (i.e., the project license). Whereas, %our preliminary study on \sihan{XX} projects shows that only  \sihan{XX}\% incompatibility issues can be fixed by recommending a new license from these candidates. 
this paper observes that %many license incompatibility issues can not be resolved only by official licenses. Moreover, 
there may exist other licenses whose copyright holders are the same with the project license, and thus can also be changed by the licensor for resolving license incompatibility issues. 

\begin{figure*}[]
    \centering
    \includegraphics[width=0.95\textwidth]{figs/overview.pdf}
    %\vspace{-2mm}
    \caption{Overview of \tool}
    \label{fig_overview}
\end{figure*}

To address the aforementioned issues, it is desirable to provide a fine-grained, scalable, and flexible solution to resolve license incompatibility issues.
To this end, we propose \tool, an %automatic
{automated} tool to fix license incompatibility issues for open source software. Specifically, given an OSS, it first extracts all licenses accompanied with the project, and organizes them with a hierarchical structure that represents the licensing scope. Then, it provides a fine-grained understanding of the regulations stated by each license, based on which \tool detects and localizes license incompatibility issues. Finally, it resolves the constraints hidden in the licensing context, and provides flexible suggestions for users to choose from. Comprehensive experiments on real-world projects demonstrate the effectiveness of \tool, with {4.09\%} \revise{false positive (FP)} rate and {0.02\%} \revise{false negative (FN)} rate for incompatibility issue localization, and {62.61\%} success rate for resolving 230 incompatible OSS on GitHub. %Furthermore, we also evaluate the impacts of license hierarchy, copyright holder, as well as license exceptions on the effectiveness of incompatibility resolution. 
%The results exhibit the importance of license incompatibility resolution and the practicability of \tool. 
%Finally, we discuss the feedback from OSS developers, the lessons learned from this work, as well as the limitations and threats to validity. 
We made all datasets and the replication package publicly available to facilitate follow-up research.



In summary, we made the following novel contributions:
\begin{itemize}
    \item We propose \tool, a fine-grained, scalable, and flexible approach to automatically interpret licenses, localize license incompatibility issues, and provide licensors with useful suggestions for resolving issues. %and made all datasets and source code publicly available. %which provides software owners with useful suggestions to reduce the legal risks induced by incorporating third-party packages.
    \item Comprehensive experiments demonstrate the effectiveness of \tool, with 4.09\% FP rate and 0.02\% FN rate for incompatibility issue localization, and 62.61\% of 230 real-world incompatible projects resolved by \tool.  
    \item We further investigate the impacts of license hierarchy,
copyright holders, and license exceptions on the effectiveness
of license incompatibility resolution. We made all the datasets and the replication package publicly available to facilitate follow-up research~\cite{ourwebsite}. 
%on Github \ling{https://github.com/anonymous123rainy/LiResolver} 
\end{itemize}
% previous studies about how to identify and fix such issues

% the motivation and the methodology of this paper

% the contribution of this paper
